% vim: ai sw=2 spell spelllang=cs

\begin{frame}{Vlákna}

  \begin{itemize}
    \item Kontext procesoru bez samostatné VM
    \begin{itemize}
      \item Sdílí s rodičovským procesem a ostatními vlákny
    \end{itemize}

    \pause

    \item Z pohledu jádra Linuxu
    \begin{itemize}
      \item Nerozlišuje
      \item Oboje je \code{struct task_struct}
      \item Oboje vzniká pomocí \code{clone}
    \end{itemize}

    \pause

    \item Použít vlákna nebo procesy?
    \begin{itemize}
      \item Velmi záleží na aplikaci
      \item Nemá jasné výhody nebo nevýhody
    \end{itemize}
  \end{itemize}

\end{frame}

\begin{frame}{Vytvoření vlákna}

  \begin{itemize}
    \item Nejčastěji používaná implementace: POSIX threads (pthreads)
    \begin{itemize}
      \item Knihovna \file{pthread} (\header{pthread.h}, \doc{man 7 pthreads})
    \end{itemize}

    \pause

    \item Vytvoření vlákna: \code{int pthread_create(pthread_t *thread, const
      pthread_attr_t *attr,}\\
      \code{void *(*start_routine) (void *), void *arg)} %stopzone
    \begin{itemize}
      \item \code{thread}: návratová hodnota -- držátko
      \item \code{attr}: atributy k vytvoření nebo \code{NULL}
      \item \code{start_routine}: funkce vlákna
      \item \code{arg}: předá se jako parametr do \code{start_routine}
      \item Návratová hodnota: \code{0} -- v pořádku, jinak číslo chyby
    \end{itemize}

    \pause

    \item Čekaní na skončení: \code{int pthread_join(pthread_t thread,
      void **retval)}
    \begin{itemize}
      \item \code{thread}: držátko z \code{pthread_create}
      \item \code{retval}: co vrátila \code{start_routine}
      \item Návratová hodnota: \code{0} -- v pořádku, jinak číslo chyby
    \end{itemize}
  \end{itemize}

\end{frame}

\begin{frame}{Úkol}

  \heading{Vytvoření vlákna}
  \begin{enumerate}
    \item Otevřete si \file{14/main.pro}
    \item Vytvořte vlákno
    \item Něco v něm vypište
    \item Vraťte z vlákna \code{0x100}
    \item V rodiči si počkejte na vlákno
    \item Vypište návratovou hodnotu vlákna
    \item Vyzkoušejte
  \end{enumerate}

\end{frame}

\begin{frame}{Ukončení vlákna}

  \begin{itemize}
    \item Vlákno lze ukončit z jiného vlákna
    \begin{itemize}
      \item Podmínka: vlákno občas vykoná jeden z ,,cancellation points''
      \item \code{sleep}, \code{open}, \code{read}, \code{close}, \code{wait},
        \dots (\doc{man pthreads})
    \end{itemize}

    \pause

    \item Ukončení: \code{int pthread_cancel(pthread_t thread)}
    \begin{itemize}
      \item \code{thread}: držátko z \code{pthread_create}
      \item Návratová hodnota: \code{0} -- v pořádku, jinak číslo chyby
      \item \code{pthread_join} dostane \code{PTHREAD_CANCELED}
    \end{itemize}
  \end{itemize}

\end{frame}

\begin{frame}{Úkol}

  \heading{Ukončení vlákna}
  \begin{enumerate}
    \item Ve vlákně spěte v cyklu
    \item V rodiči spěte 2 vteřiny
    \item Potom ukončete vlákno
    \item A počkejte si na vlákno
    \item Vypište návratovou hodnotu vlákna
    \item Vyzkoušejte
  \end{enumerate}

\end{frame}
